\documentclass[12pt,a4paper]{article}
\usepackage[margin=2cm]{geometry}
\usepackage{graphicx,xcolor,tabularx,booktabs,array}
\usepackage[T1]{fontenc}
\usepackage[utf8]{inputenc}
\usepackage{caption}
\captionsetup[table]{skip=6pt}
\title{Nahuales Mayas: Glifo y Colores (glifos nuevos)}
\date{}
\begin{document}
\maketitle
\begin{table}[h!]
\centering
\renewcommand{\arraystretch}{1.3}
\begin{tabularx}{\textwidth}{|>{\raggedright\arraybackslash}p{3cm}|c|X|}
\hline
\textbf{Nahual} & \textbf{Glifo} & \textbf{Colores} \\ \hline
B'ATZ' & \includegraphics[height=1.8cm]{{glyphs_new/glyph_1.png}} & Rojo, Anaranjado \\ \hline
E & \includegraphics[height=1.8cm]{{glyphs_new/glyph_5.png}} & Blanco, Celeste \\ \hline
AJ & 
\includegraphics[height=1.8cm]{{glyphs_new/glyph_8.png}} & Caf\'e, blanco \\ \hline
I'X & \includegraphics[height=1.8cm]{{glyphs_new/glyph_4.png}} & Amarillo, caf\'e \\ \hline
TZ'IKIN & 
\includegraphics[height=1.8cm]{{glyphs_new/glyph_18.png}} & Azul, celeste, blanco \\ \hline
AJMAQ & \includegraphics[height=1.8cm]{{glyphs_new/glyph_10.png}} & Gris, negro, blanco \\ \hline
NO'J & 
\includegraphics[height=1.8cm]{{glyphs_new/glyph_20.png}} & Azul, celeste \\ \hline
TIJAX & \includegraphics[height=1.8cm]{{glyphs_new/glyph_6.png}} & Blanco, rojo, negro \\ \hline
AQ'AB'AL & \includegraphics[height=1.8cm]{{glyphs_new/glyph_13.png}} & Rojo, azul marino, anaranjado \\ \hline
K'AT & 
\includegraphics[height=1.8cm]{{glyphs_new/glyph_17.png}} & Caf\'e, amarillo, beige \\ \hline
KAN & \includegraphics[height=1.8cm]{{glyphs_new/glyph_15.png}} & Verde, rojo \\ \hline
KAME & \includegraphics[height=1.8cm]{{glyphs_new/glyph_14.png}} & Amarillo, anaranjado, negro \\ \hline
KEJ & 
\includegraphics[height=1.8cm]{{glyphs_new/glyph_3.png}} & Beige, amarillo \\ \hline
Q'ANIL & \includegraphics[height=1.8cm]{{glyphs_new/glyph_7.png}} & Gris, blanco, verde \\ \hline
TOJ & \includegraphics[height=1.8cm]{{glyphs_new/glyph_11.png}} & Rojo, blanco \\ \hline
TZ'I' & \includegraphics[height=1.8cm]{{glyphs_new/glyph_2.png}} & Amarillo, blanco, beige \\ \hline
\end{tabularx}
\end{table}

\small Fuente: glifos reemplazados a partir de la imagen de 20 glifos proporcionada por el usuario. 
\end{document}
